% GENERATED FILE. Edit canonical lessons, then run npm run generate:books.
% canonical-source-hash: fnv1a64:3586a67723f60d5e
% canonical-lessons: FA-C504-dusti, FA-C504-dust-shodan, FA-C504-rabete, FA-C504-gharibe, FA-C504-hamdigar

\chapter{Friends}
\label{ch:fa-friends}
\hlchaptermodality{504}

\begin{chapteropening}
\textbf{By the end of this chapter:} \emph{I can talk about friendship.}

Complete the last lesson of chapter 504: I can talk about friendship.
\end{chapteropening}

\section[\texorpdfstring{\fa{دوستی} (dusti)}{دوستی (dusti)}]{\fa{دوستی} (dusti) — friendship}

\label{lesson:FA-C504-dusti}



\begin{quote}
Before the new one: say the Persian for spoiled, then the Persian for hide-and-seek.
\end{quote}

\subsection*{You'll want to know: \fa{دوستی}}
\textbf{\fa{دوستی}} (\emph{dusti}) — \textquotedblleft{}friendship\textquotedblright{}.

\begin{quote}
\fa{دوستی} \fa{ما} \fa{از} \fa{دبیرستان} \fa{شروع} \fa{شد}.
\end{quote}

\begin{quote}
\emph{dusti-ye mâ az dabirestân shoru' shod.}
\end{quote}

\begin{quote}
Our friendship began in high school.
\end{quote}

\emph{dust} 'friend' + \emph{-i}. As a modifier, 'friendly': \emph{mosâbeqe-ye dusti} 'a friendly match'.

\begin{cousinweb}
\emph{dusti} is built on native \emph{dust}, 'friend', which is usually traced to Old Persian \emph{dauštā}, 'beloved', with the noun ending \emph{-i}. EtymDB lists \emph{dust} without a source.
\end{cousinweb}

\begin{grammarlens}[title={What you've built}]
One more word for telling, explaining and describing at B1.
\end{grammarlens}

\subsection*{Your turn}
\begin{itemize}
\raggedright
  \item \emph{Say it:} \emph{dusti}
  \item \emph{Say it:} \emph{dusti}, once more
  \item \emph{Say it:} a sentence of your own with \emph{dusti}
  \item \emph{Recall:} say the Persian for spoiled, then the Persian for hide-and-seek, then say \emph{dusti} again
\end{itemize}

\begin{tcolorbox}[breakable,colback=teal!4,colframe=teal!35!black,title={Before you move on}]
What does \fa{دوستی} mean? (\textquotedblleft{}Friendship\textquotedblright{}.) Say it once more.
\end{tcolorbox}
\section[\texorpdfstring{\fa{دوست} \fa{شدن} (dust shodan)}{دوست شدن (dust shodan)}]{\fa{دوست} \fa{شدن} (dust shodan) — to make friends}

\label{lesson:FA-C504-dust-shodan}



\begin{quote}
Before the new one: say the Persian for hide-and-seek, then the Persian for friendship.
\end{quote}

\subsection*{You'll want to know: \fa{دوست} \fa{شدن}}
\textbf{\fa{دوست} \fa{شدن}} (\emph{dust shodan}) — \textquotedblleft{}to make friends\textquotedblright{}.

\begin{quote}
\fa{او} \fa{در} \fa{مدرسه} \fa{جدید} \fa{زود} \fa{با} \fa{همه} \fa{دوست} \fa{شد}.
\end{quote}

\begin{quote}
\emph{u dar madrese-ye jadid zud bâ hame dust shod.}
\end{quote}

\begin{quote}
At the new school he quickly made friends with everyone.
\end{quote}

\emph{bâ kasi dust shodan}; \emph{dust} 'friend' + \emph{shodan} 'to become'.

\begin{cousinweb}
\emph{dust shodan} is built on two inherited words: \emph{dust}, 'friend', usually traced to Old Persian \emph{dauštā}, 'beloved', and \emph{shodan}, 'to become'.
\end{cousinweb}

\begin{grammarlens}[title={What you've built}]
One more word for telling, explaining and describing at B1.
\end{grammarlens}

\subsection*{Your turn}
\begin{itemize}
\raggedright
  \item \emph{Say it:} \emph{dust shodan}
  \item \emph{Say it:} \emph{dust shodan}, once more
  \item \emph{Say it:} a sentence of your own with \emph{dust shodan}
  \item \emph{Recall:} say the Persian for hide-and-seek, then the Persian for friendship, then say \emph{dust shodan} again
\end{itemize}

\begin{tcolorbox}[breakable,colback=teal!4,colframe=teal!35!black,title={Before you move on}]
What does \fa{دوست} \fa{شدن} mean? (\textquotedblleft{}To make friends\textquotedblright{}.) Say it once more.
\end{tcolorbox}
\section[\texorpdfstring{\fa{رابطه} (râbete)}{رابطه (râbete)}]{\fa{رابطه} (râbete) — a relationship; a connection}

\label{lesson:FA-C504-rabete}



\begin{quote}
Before the new one: say the Persian for friendship, then the Persian for to make friends.
\end{quote}

\subsection*{You'll want to know: \fa{رابطه}}
\textbf{\fa{رابطه}} (\emph{râbete}) — \textquotedblleft{}a relationship; a connection\textquotedblright{}.

\begin{quote}
\fa{رابطه} \fa{او} \fa{با} \fa{پدرش} \fa{خیلی} \fa{خوب} \fa{است}.
\end{quote}

\begin{quote}
\emph{râbete-ye u bâ pedarash kheyli khub ast.}
\end{quote}

\begin{quote}
Her relationship with her father is very good.
\end{quote}

Plural \emph{ravâbet} (Arabic) in formal Persian; \emph{râbetehâ} in speech.

\begin{cousinweb}
\emph{râbete} comes from Arabic \emph{rābiṭa}, 'bond, tie' (EtymDB).
\end{cousinweb}

\begin{grammarlens}[title={What you've built}]
One more word for telling, explaining and describing at B1.
\end{grammarlens}

\subsection*{Your turn}
\begin{itemize}
\raggedright
  \item \emph{Say it:} \emph{râbete}
  \item \emph{Say it:} \emph{râbete}, once more
  \item \emph{Say it:} a sentence of your own with \emph{râbete}
  \item \emph{Recall:} say the Persian for friendship, then the Persian for to make friends, then say \emph{râbete} again
\end{itemize}

\begin{tcolorbox}[breakable,colback=teal!4,colframe=teal!35!black,title={Before you move on}]
What does \fa{رابطه} mean? (\textquotedblleft{}A relationship; a connection\textquotedblright{}.) Say it once more.
\end{tcolorbox}
\section[\texorpdfstring{\fa{غریبه} (gharibe)}{غریبه (gharibe)}]{\fa{غریبه} (gharibe) — a stranger}

\label{lesson:FA-C504-gharibe}



\begin{quote}
Before the new one: say the Persian for to make friends, then the Persian for a relationship.
\end{quote}

\subsection*{You'll want to know: \fa{غریبه}}
\textbf{\fa{غریبه}} (\emph{gharibe}) — \textquotedblleft{}a stranger\textquotedblright{}.

\begin{quote}
\fa{بچه} \fa{ها} \fa{با} \fa{غریبه} \fa{ها} \fa{حرف} \fa{نمیزنند}.
\end{quote}

\begin{quote}
\emph{bachchehâ bâ gharibehâ harf nemizanand.}
\end{quote}

\begin{quote}
The children don't talk to strangers.
\end{quote}

Built on \emph{gharib} 'strange, foreign' with \emph{-e}. Plural \emph{gharibehâ}.

\begin{cousinweb}
\emph{gharibe} is built in Persian on \emph{gharib}, 'strange', which was borrowed from Arabic \emph{ġarīb}, 'outsider, foreigner' (EtymDB), with the native ending \emph{-e}.
\end{cousinweb}

\begin{grammarlens}[title={What you've built}]
One more word for telling, explaining and describing at B1.
\end{grammarlens}

\subsection*{Your turn}
\begin{itemize}
\raggedright
  \item \emph{Say it:} \emph{gharibe}
  \item \emph{Say it:} \emph{gharibe}, once more
  \item \emph{Say it:} a sentence of your own with \emph{gharibe}
  \item \emph{Recall:} say the Persian for to make friends, then the Persian for a relationship, then say \emph{gharibe} again
\end{itemize}

\begin{tcolorbox}[breakable,colback=teal!4,colframe=teal!35!black,title={Before you move on}]
What does \fa{غریبه} mean? (\textquotedblleft{}A stranger\textquotedblright{}.) Say it once more.
\end{tcolorbox}
\section[\texorpdfstring{\fa{همدیگر} (hamdigar)}{همدیگر (hamdigar)}]{\fa{همدیگر} (hamdigar) — each other}

\label{lesson:FA-C504-hamdigar}



\begin{quote}
Before the new one: say the Persian for a relationship, then the Persian for a stranger.
\end{quote}

\subsection*{You'll want to know: \fa{همدیگر}}
\textbf{\fa{همدیگر}} (\emph{hamdigar}) — \textquotedblleft{}each other\textquotedblright{}.

\begin{quote}
\fa{آنها} \fa{از} \fa{کودکی} \fa{همدیگر} \fa{را} \fa{میشناسند}.
\end{quote}

\begin{quote}
\emph{ânhâ az kudaki hamdigar râ mishenâsand.}
\end{quote}

\begin{quote}
They've known each other since childhood.
\end{quote}

\emph{ham} 'together' + \emph{digar} 'other'. Formal Persian also says \emph{yekdigar}.

\begin{cousinweb}
\emph{hamdigar} is built on two native Persian words: \emph{ham}, 'together, same', which EtymDB traces to Middle Persian (EtymDB), and \emph{digar}, 'other'.
\end{cousinweb}

\begin{grammarlens}[title={What you've built}]
One more word for telling, explaining and describing at B1.
\end{grammarlens}

\subsection*{Your turn}
\begin{itemize}
\raggedright
  \item \emph{Say it:} \emph{hamdigar}
  \item \emph{Say it:} \emph{hamdigar}, once more
  \item \emph{Say it:} a sentence of your own with \emph{hamdigar}
  \item \emph{Recall:} say the Persian for a relationship, then the Persian for a stranger, then say \emph{hamdigar} again
\end{itemize}

\begin{tcolorbox}[breakable,colback=teal!4,colframe=teal!35!black,title={Before you move on}]
What does \fa{همدیگر} mean? (\textquotedblleft{}Each other\textquotedblright{}.) Say it once more.
\end{tcolorbox}
